\section*{Chapitre \ref{ch:ml:the-machine-learning-team} --- L'équipe d'apprentissage automatique}

\begin{solution}{exo:ml:the-machine-learning-team:1}
Les chercheurs sont \omterm{def:ml:the-machine-learning-team:roles}{propriétaires du modèle} et de ses variables (le signal, l'ajustement, les éléments de validation)~; les \omterm{def:ml:the-machine-learning-team:roles}{ingénieurs en apprentissage automatique} (ML engineers) ont en charge le
service qui l'exécute (variables en ligne, modèle compilé, moniteurs)~; l'équipe plateforme possède tout ce que partagent tous les modèles (\omterm{def:ml:data-and-feature-stores:store}{feature store},
entraînement, registre, service, validateur de transmission).
\end{solution}

\begin{solution}{exo:ml:the-machine-learning-team:2}
Parce que quelqu'un doit répondre à une \omterm{def:ml:monitoring-and-retraining:performance-alarm}{alarme de performance}, décider du réentraînement ou du retrait et expliquer le modèle aux validateurs et
aux régulateurs. Une équipe propriétaire dilue les trois~: l'alerte part vers une liste, chaque membre suppose qu'un autre agit, et le modèle tourne
sans surveillance.
\end{solution}

\begin{solution}{exo:ml:the-machine-learning-team:3}
$L=\lambda W=0.8\times5.05=4.04$ modèles, contre 4.00 pour la chaîne (théorie~: $\rho/(1-\rho)=4$ avec $\rho=0.8$).
\end{solution}

\begin{solution}{exo:ml:the-machine-learning-team:4}
Une fenêtre de 5\,000 est une définition valide (une entrée connue, une agrégation connue, une fenêtre positive), et les calculs hors ligne et en ligne
utilisent tous deux cette fenêtre, si bien qu'ils concordent entre eux. Seule la comparaison des valeurs calculées de la variable avec les valeurs propres à la recherche
sur la session de référence montre que la variable n'est pas celle sur laquelle le modèle a été entraîné.
\end{solution}

\begin{solution}{exo:ml:the-machine-learning-team:5}
Dans une file M/M/1, l'attente moyenne avant le service vaut $\rho/(1-\rho)$ durées de service~: 2.8 à 74~\% et 9 à 90~\%, soit un peu plus de trois
fois plus longue. Le poste du chapitre compte deux serveurs et des reprises, si bien que les chiffres diffèrent, mais la forme (une attente qui explose à l'approche
de la saturation) est la même, comme le montre la \cref{fig:ml:team:rate}.
\end{solution}

\begin{solution}{exo:ml:the-machine-learning-team:6}
La validation élimine 7.9~\% des modèles, la revue en élimine 41~\%, et le temps se perd dans les files d'attente et les reprises chez les ingénieurs, non dans les
verdicts de la validation. Assouplir les validateurs laisserait passer davantage de mauvais modèles vers le trading simulé (paper trading) et la production sans raccourcir beaucoup le
délai~; supprimer la réécriture (la boucle de reprises entre les ingénieurs et les validateurs) est ce qui le raccourcit.
\end{solution}

\begin{solution}{exo:ml:the-machine-learning-team:7}
\texttt{ml\_team.third\_engineer}~: avec trois ingénieurs, le taux d'utilisation des ingénieurs tombe à 50~\% et le délai médian à 16.6
semaines (moyennes sur trois graines), contre 22.3 avec deux ingénieurs et 11.1 avec le paquet, sur les mêmes graines. Les effectifs suppriment
une partie de l'attente dans les files, mais aucune des reprises~; le paquet supprime les deux et libère le temps des ingénieurs au lieu de l'alourdir. Le
paquet d'abord~; un troisième ingénieur seulement si la file revient à mesure que la recherche grandit.
\end{solution}

\begin{solution}{exo:ml:the-machine-learning-team:8}
Artefact~: le vecteur des coefficients et l'ordonnée à l'origine, avec leur empreinte de contenu. Variables~: cinq définitions avec unités, sources et la
standardisation ajustée sur les 60 jours d'entraînement (moyennes et écarts-types faisant partie de l'artefact). Vecteurs de variables~: entrées
brutes et valeurs standardisées sur les jours de référence. Vecteurs de test~: 200 lignes de variables avec prédictions à $10^{-12}$. Latence~: négligeable,
mais énoncée. Surveillance~: les cinq moniteurs du chapitre~27 avec leurs seuils (0.115, 0.115, 0.089, 0.07 et les paramètres $k$ et
$h$ du CUSUM), une fausse alerte par mois, avec les propriétaires. \omterm{def:ml:interpretability-and-model-governance:card}{Fiche de modèle} (model card)~: finalité, cible (rendement du lendemain), horizon, seuil de trading, limites
(linéaire, sensible aux unités). Contrôles~: comme dans le chapitre, plus un contrôle de plage sur chaque entrée brute, qui aurait détecté le
mardi du chapitre~27 avant que le modèle ne le voie.
\end{solution}

\begin{solution}{pb:ml:the-machine-learning-team:1}
\textbf{Partie I.}
\begin{enumerate}
\item Les chercheurs, les \omterm{def:ml:the-machine-learning-team:roles}{ingénieurs en apprentissage automatique} (ML engineers), les ingénieurs de plateforme, le validateur de modèles, le trader ou le gérant de portefeuille, et le
\omterm{def:ml:the-machine-learning-team:roles}{propriétaire du modèle} (un rôle tenu par l'un d'eux)~; Kreuzberger et coauteurs en listent sept autour du MLOps.
\item La performance, la surveillance, le réentraînement et le retrait, et l'explication du modèle à la validation, au risque et aux régulateurs.
\item Tout ce que partagent tous les modèles (\omterm{def:ml:data-and-feature-stores:store}{feature store}, entraînement, suivi, registre, service, moniteurs, validateur de transmission)~; non les
modèles ni leurs variables.
\item Parce que la répartition détermine ce qui circule entre les équipes~: une réécriture crée des files et des reprises, une base de code commune avec un
paquet vérifié n'en crée pas~; les modèles sont plus difficiles à modulariser que les logiciels (Amershi et coauteurs).
\end{enumerate}
\textbf{Partie II.}
\begin{enumerate}[resume]
\item La forêt et son empreinte, 16 \omterm{def:ml:data-and-feature-stores:store}{définitions de variables}, les valeurs des variables de la recherche sur une session de référence, 200 vecteurs de test à
$10^{-12}$, un budget de 5 microsecondes avec 2.6 mesurées, quatre moniteurs avec propriétaires, une \omterm{def:ml:interpretability-and-model-governance:card}{fiche de modèle} de quinze champs.
\item Schéma, empreinte de l'artefact, spécification des variables, parité des variables, vecteurs de variables, vecteurs de test, latence, surveillance, \omterm{def:ml:interpretability-and-model-governance:card}{fiche de modèle}.
\item Le \Cref{tab:ml:team:defects}~: artefact reconstruit, empreinte et vecteurs de test~; entrée inconnue et variable manquante, la spécification~;
fenêtre en millisecondes, vecteurs de variables~; sorties d'un autre modèle, vecteurs de test~; latence de la bibliothèque, latence~; moniteur sans propriétaire,
surveillance~; fiche sans limites, la fiche.
\item Une variable erronée peut être calculée à l'identique hors ligne et en ligne~; seule une comparaison avec les valeurs de la recherche la détecte.
\end{enumerate}
\textbf{Partie III.}
\begin{enumerate}[resume]
\item Arrivées poissonniennes à 0.35 par semaine~; revue (élimine 35~\%, 15~\% renvoyés à la recherche), réimplémentation par deux ingénieurs (30~\% renvoyés
pour clarification), validation (élimine 10~\%, 25~\% renvoyés aux ingénieurs), quatre semaines de trading simulé (élimine 20~\%), un
canari de deux semaines (élimine 5~\%).
\item Face à la file M/M/1~: 5.05 semaines simulées contre 5, 4.00 modèles dans le système selon la chaîne du livre~4, 4.04 par la loi de Little~;
sur l'ensemble du pipeline, 4.87 contre 4.90.
\item La plupart des pertes surviennent à la revue (41~\%)~; le temps part dans les boucles de reprise et dans la file d'attente chez les ingénieurs, occupés 74~\% du temps, chaque modèle
passant par elles 1.79 fois.
\item Avec la réécriture, le délai médian passe d'environ 16 semaines pour 0.2 modèle par semaine à 70 semaines pour 0.44~; avec le paquet, il reste voisin de 11 semaines.
\end{enumerate}
\textbf{Partie IV.}
\begin{enumerate}[resume]
\item \emph{Qui est \omterm{def:ml:the-machine-learning-team:roles}{propriétaire du modèle}~?} Avec une réécriture à la transmission, le modèle médian atteint la production 20.4 semaines après la recherche
(90\up{e}~centile 42.9), et 35~\% y parviennent~: 41~\% sont perdus à la revue, 5~\% à la réimplémentation, 8~\% à la validation, 9~\% en
trading simulé (paper trading) et 2~\% au canari. Un paquet validé à la place de la réécriture ramène le délai médian à 11.1 semaines, soit un gain de 9.3 semaines.
\item Parce que les portes de contrôle jugent les modèles, et que les modèles n'ont pas changé~; seul le temps passé entre les portes devrait changer.
\item Une semaine de temps de recherche par modèle~; elle fait gagner 9.3 semaines de délai et ramène la charge des ingénieurs de 74~\% à 12~\%.
\item Non~: 16.6 semaines contre 11.1, puisqu'il raccourcit la file mais conserve les reprises.
\item Délai médian et délai au 90\up{e}~centile, pertes par porte, taux d'utilisation par équipe, taux de reprise par transmission.
\item Le chercheur qui l'a construit ou le responsable du desk qui le trade, nommé dans le \omterm{def:ml:experiment-tracking-and-reproducibility:registry}{registre de modèles}, et non l'équipe de plateforme.
\item Le modèle retourne à son auteur avec le détail du contrôle échoué, et rien n'atteint la production tant qu'il n'a pas réussi.
\item Une base de code commune à la recherche et à la production, et un paquet qu'une machine peut vérifier.
\end{enumerate}
\end{solution}

\begin{solution}{iq:ml:the-machine-learning-team:1}
L'ingénieur construit et fait tourner le chemin de production~: variables en ligne identiques aux variables hors ligne, un modèle dans son budget de latence, déploiement,
surveillance, chaînes de réentraînement et retours arrière. Le chercheur trouve et valide le signal~; l'ingénieur le fait tourner sans
son auteur.
\iqlookfor{enjeux de production~: parité, latence, surveillance, exploitation.}
\end{solution}

\begin{solution}{iq:ml:the-machine-learning-team:2}
L'artefact et son empreinte, les \omterm{def:ml:data-and-feature-stores:store}{définitions de variables} dans le format commun avec des valeurs de référence, des vecteurs de test d'entrées et
de sorties, un budget de latence et sa mesure, la spécification de surveillance, une \omterm{def:ml:interpretability-and-model-governance:card}{fiche de modèle} et un propriétaire, le tout vérifiable automatiquement.
\iqlookfor{un paquet vérifiable, pas un notebook.}
\end{solution}

\begin{solution}{iq:ml:the-machine-learning-team:3}
Comparer les valeurs hors ligne et en ligne aux mêmes instants de décision sur un échantillon (parité) et aux valeurs de référence de la recherche~;
repérer la première variable et la première entrée qui divergent. Empêcher cela avec une seule \omterm{def:ml:data-and-feature-stores:store}{définition de variable} exécutée par les deux chemins (\omterm{def:ml:data-and-feature-stores:store}{feature store}),
contrôles de parité en intégration continue et en production, et vecteurs de test des variables à chaque transmission.
\iqlookfor{tests de parité et définition unique.}
\end{solution}

\begin{solution}{iq:ml:the-machine-learning-team:4}
Enregistrer les horodatages de chaque modèle à chaque étape, et mesurer l'attente par rapport au temps de travail, les boucles de reprise et le taux d'utilisation de chaque
équipe~; modéliser le pipeline comme des files pour tester les changements. Le plus souvent, le temps se perd dans les files devant une équipe saturée et dans les reprises causées par
la transmission.
\iqlookfor{mesure par étape, files et reprises.}
\end{solution}

\begin{solution}{iq:ml:the-machine-learning-team:5}
La plateforme~: stockage et calcul des variables, infrastructure d'entraînement, suivi et registre, compilation et service,
surveillance, contrôles de transmission. L'équipe du modèle~: le modèle, ses variables et ses étiquettes, ses éléments de validation, ses
seuils de surveillance et sa propriété.
\iqlookfor{infrastructure commune face aux choix propres à chaque modèle.}
\end{solution}

\begin{solution}{iq:ml:the-machine-learning-team:6}
Le taux d'utilisation de chaque équipe et les taux de reprise aujourd'hui~: doubler les arrivées vers une équipe déjà occupée aux trois quarts du temps
fait exploser les délais bien avant que l'équipe ne soit à court d'heures. Décider ensuite de ce qu'il faut automatiser (contrôles, empaquetage) avant d'embaucher.
\iqlookfor{taux d'utilisation et files d'attente à l'approche de la saturation.}
\end{solution}
